%──────────────────────────────────────────────────────────────────
\subsection{Summary of analytic predictions and comparison with
            numerics}
\label{subsec:predictions-summary}
%──────────────────────────────────────────────────────────────────

We collect the analytic predictions of this section, organised by the
level of theoretical support, and contrast them with the numerical
results obtained from the cloning data of
Section~\ref{sec:binder-analysis}.

\begin{table}[h]
\centering
\renewcommand{\arraystretch}{1.3}
\begin{tabular}{p{0.36\linewidth}p{0.27\linewidth}p{0.27\linewidth}}
\toprule
\textbf{Quantity} & \textbf{Case~A} & \textbf{Case~B}\\
\midrule
AZ class & $D$ & $\mathrm{DIII}$\\
NLSM target & $SO(2N)/U(N)$ & $SO(N)$\\
Self-duality & yes ($\alpha\leftrightarrow\gamma$) & no\\
Critical line position & $\lambda_c^{\rm A}=\tfrac12\ \forall\zeta$ & varies with $\zeta$\\
$\lambda_c$ at small $\zeta$ & $\tfrac12$ (pinned) & $\propto\sqrt{\zeta}$\\
$\lambda_c$ at $\zeta=1$ & $\tfrac12$ & $\to\tfrac12$ (BR/Carollo)\\
Crossover form (Case~B) & --- & $\lambda_c=\sqrt{\zeta}/(1+\sqrt{\zeta})$\\
UV scaling dim.\ of $\mathcal V_\zeta$ & $1$ & $1$\\
NLSM correlation exp.\ $\nu$ & $1$ (Ising at $\theta=\pi$) & $2$ (DIII multicritical)\\
$\pi_2$ of target / topological term & $\mathbb Z$, $\theta\in\{0,\pi\}$
                                     & $\mathbb Z_2$ ($N\geq 3$)\\
Critical CFT (at $\theta=\pi$) & Ising ($c=\tfrac12$) & free-fermion Gaussian\\
\bottomrule
\end{tabular}
\caption{Analytic predictions for QJ--PPS Cases~A and~B.  Case~A is
protected by self-duality and has $\zeta$-independent universality.
Case~B has no protection; the matched-NLSM framework of
Section~\ref{subsec:matching} predicts a $\sqrt{\zeta}$ small-$\zeta$
asymptote, with a Born-rule endpoint at $\lambda_c(1)=1/2$, smoothly
interpolated by the closed-form $\lambda_c=\sqrt{\zeta}/(1+\sqrt{\zeta})$.}
\label{tab:caseAB-summary}
\end{table}


\paragraph{Comparison with cloning data (Case~B).}
The principal prediction of Section~\ref{subsec:matching}---that the
phase boundary satisfies $\alpha_c/w\propto\sqrt{\zeta}$ with
order-unity prefactor and Born-rule limit $\lambda_c\to 1/2$ at
$\zeta\to 1$---is tested in
Section~\ref{sec:binder-analysis} via global finite-size scaling
collapse on the cloning data at $L\in\{32,48,64,96,128,192,256\}$.
On the cleanest data range $\zeta\in[0.03,0.85]$ (excluding the
finite-size-limited $\zeta=1$ corner and the asymptotic-FSS-marginal
$\zeta=0.02$ corner), a free power-law fit
$\alpha_c/w=C\,\zeta^{\phi}$ gives
\begin{equation}\label{eq:phi-C-fit}
  \phi^{\rm fit}\,=\,0.56\pm 0.05,
  \qquad
  C^{\rm fit}\,=\,1.02\pm 0.10,
  \qquad
  \chi^{2}/\mathrm{dof}\,=\,3.8.
\end{equation}
The prediction $\phi=1/2$ of Eq.~\eqref{eq:caseB-sqrtz} is consistent
within $1.3\sigma$; the alternative $\phi=1$, corresponding to a naive
$\xi_\lambda^{\rm cross}\propto\lambda^{-1}$ matching length, is
excluded at $9\sigma$.  The fitted prefactor $C\approx 1$ is
consistent with the order-unity expectation; we emphasise that $C$ is
non-universal and that no precision target is meaningful at this
level.

\paragraph{Logarithmic corrections in the log phase.}
The one-loop DIII flow~\eqref{eq:caseB-beta-R1} predicts a specific
form for the logarithmic corrections in the weakly-monitored phase of
Case~B.  Integrating Eq.~\eqref{eq:caseB-beta-R1} gives
\begin{equation}\label{eq:running-coupling}
  g_R(L) \;=\; \frac{g_B}{1 + (g_B/8\pi)\ln(L/L_{\rm match})},
\end{equation}
where $L_{\rm match}$ is the matching scale (either $a$ or
$\xi_\lambda^{\rm cross}$ depending on the regime).  This predicts a
sub-leading correction to the entanglement entropy of the form
$S(L) = (c_{\rm eff}/6)\ln L + b\ln\ln L+\ldots$ with $b$ a universal
coefficient set by the $\beta$-function.  The drift of $c_{\rm eff}(L)$
with $L$ at fixed $\zeta$ is in principle a direct signature of the
DIII flow, although the available cloning data is not precise enough
to extract $b$ reliably at the $L\leq 256$ system sizes accessible
here.

\paragraph{R\'enyi-ratio drift.}
The weak-coupling fixed point of the DIII NLSM is the free-fermion
Gaussian CFT, with R\'enyi ratios
$c_n/c_1 = \tfrac{1}{2}(1+1/n)$, i.e.\ $c_2/c_1=0.75$,
$c_3/c_1\approx 0.667$.  The cloning data of
Section~\ref{sec:cloning} measure $c_2/c_1$ at $\lambda=0.05$ deep in
the log phase across $L\in\{64,96,128,192,256,384\}$ and find values
in the range $[0.75,0.80]$ at $\zeta=0.5$ and $[0.81,0.84]$ at
$\zeta=1$.  The framework predicts these ratios should approach the
Gaussian values from above as $L\to\infty$, but the available data
shows a non-monotonic drift at $\zeta=0.5$ and approximate stability
at $\zeta=1$.  The clean drift signal is masked by reduced trajectory
counts at the larger system sizes ($N_c$ falls from $200$ at $L=64$
to $20$ at $L=384$), and the test is therefore inconclusive at the
present statistics.

\paragraph{Two-parameter FSS in $(\lambda L^{1/2},\zeta L)$.}
At the multicritical fixed point, the joint scaling form
\begin{equation}\label{eq:two-param-FSS}
  B_L(\lambda,\zeta) \;=\; \mathcal B\!\bigl(\lambda L^{1/y_\lambda},\,\zeta L^{y_\zeta}\bigr)
  \;=\; \mathcal B\!\bigl(\lambda L^{1/2},\,\zeta L\bigr)
\end{equation}
should describe the Binder cumulant up to logarithmic corrections in
the small-$\zeta$ regime.  A direct two-parameter collapse on the
cloning data shows scatter at the order-of-magnitude level at fixed
$(\lambda L^{1/2},\zeta L)$, consistent with finite-size corrections
expected at $\zeta L\gtrsim 5$ (the marginal asymptotic-FSS threshold
for the available $L$ values).  The single-parameter scaling
$B_L = F((\lambda-\lambda_c(\zeta))L^{1/\nu})$ at fixed $\zeta$ gives
clean collapses across all five $L$ values, supporting the existence
of a $\zeta$-dependent critical line of class-DIII type.

\paragraph{The empirical $\nu(\zeta)$ pattern.}
Global FSS extracts $(\lambda_c,\nu)$ jointly at each $\zeta$ by
collapsing the five $L$-curves; the resulting $\nu$ values scatter
between $1.7$ and $3.0$, with weighted mean $\bar\nu=2.23\pm 0.04$.
On the clean middle range $\zeta\in[0.1,0.7]$ the values cluster
within $1\sigma$ of $\nu=2$, consistent with the class-DIII
multicritical prediction.  The end points $\zeta=0.85$ and $\zeta=1$
pull $\bar\nu$ upward to $\sim 2.7$--$3.0$; these are precisely the
regions where finite-size systematics (low $N_c$ at $L=256$ and
limited $\lambda$ coverage) dominate, and we therefore treat $\nu=2$
as a literature input that is consistent with our data in the
asymptotic-FSS regime rather than as an independent numerical
determination.

\paragraph{Summary of theoretical robustness.}
The principal robust prediction of this chapter is the small-$\zeta$
scaling law $\lambda_c\propto\sqrt{\zeta}$ for Case~B, derived from
two field-theory inputs ($\Delta_\zeta^{\rm UV}=1$ from
cross-Choi, microscopically verified; $\nu=2$ from class-DIII
multicritical NLSM, taken from the disordered-fermion literature),
combined via the matched-RG argument of
Section~\ref{subsec:matching}.  This prediction is confirmed by the
cloning numerics within $1.3\sigma$ on a $9$-point dataset spanning
two decades of $\zeta$.  The non-universal prefactor $C\approx 1$ is
consistent with order-unity matching; precision computation of $C$
would require a one-loop NLSM calculation that is not pursued here.
The Born-rule limit $\lambda_c\to 1/2$ at $\zeta=1$ is consistent
with Carollo~\cite{Carollo2018PRA}.  These two limits are smoothly
joined by the parameter-free closed-form interpolation
$\lambda_c(\zeta) = \sqrt{\zeta}/(1+\sqrt{\zeta})$.
